\documentclass[10pt]{article}
\usepackage[T1]{fontenc}
\usepackage[utf8]{inputenc}
\usepackage[french]{babel}
\usepackage{amsmath,amssymb,mathrsfs,mathpazo}
\usepackage{geometry}
\geometry{a4paper,margin=2cm}
\usepackage{graphicx,array,tabularx,tabularray,colortbl,multirow,multicol,booktabs,enumitem,pifont,ulem,fancybox,diagbox,cancel}
\usepackage{pstricks,pstricks-add,pst-plot,pst-tree,pst-node,pst-text,pst-func,pst-math,pst-bspline,pst-3dplot,pst-eucl,pst-solides3d}
\let\pscalculate\undefined
\usepackage{pst-bezier}
\usepackage{tikz,pgfplots,tkz-tab}
\usetikzlibrary{arrows,arrows.meta,patterns,calc,positioning,trees,decorations.markings,shapes.geometric,shapes.misc,shapes.arrows,intersections,decorations.pathmorphing,decorations.pathreplacing,backgrounds,shadows,babel}
\pgfplotsset{compat=1.18}
\usepgfplotslibrary{fillbetween}
\usepackage[np]{numprint}
\usepackage{listings,makecell,pas-tableur,twoopt}
\let\barre\undefined
\usepackage[pstricks]{tablvar}
\definecolor{darkgreen}{RGB}{0,100,0}
\providecommand{\red}{\color{red}}
\providecommand{\blue}{\color{blue}}
\providecommand{\green}{\color{green}}
\providecommand{\black}{\color{black}}
\providecommand{\white}{\color{white}}
\providecommand{\gray}{\color{gray}}
\providecommand{\cyan}{\color{cyan}}
\providecommand{\magenta}{\color{magenta}}
\providecommand{\yellow}{\color{yellow}}
\providecommand{\violet}{\color{violet}}
\providecommand{\orange}{\color{orange}}
\providecommand{\euro}{\texteuro}
\providecommand{\EUR}{\texteuro}
\providecommand{\ieme}{\textsuperscript{e}}
\providecommand{\degres}{\ensuremath{^\circ}}
\providecommand{\degree}{\ensuremath{^\circ}}
\providecommand{\vv}[1]{\overrightarrow{#1}}
\providecommand{\Coord}[1]{#1}
\providecommand{\decosix}{}
\providecommand{\starredbullet}{\ensuremath{\bullet}}
\pagestyle{empty}
\setlength{\parindent}{0pt}
\usepackage{adjustbox}\DeclareUnicodeCharacter{2236}{:}\DeclareUnicodeCharacter{2190}{\ensuremath{\leftarrow}}
\providecommand{\R}{\mathbb{R}}
\providecommand{\N}{\mathbb{N}}
\providecommand{\Z}{\mathbb{Z}}
\providecommand{\Q}{\mathbb{Q}}
\providecommand{\D}{\mathbb{D}}
\providecommand{\C}{\mathbb{C}}
\providecommand{\e}{{\mathrm{e}}}
\providecommand{\vect}[1]{\overrightarrow{#1}}
\providecommand{\vectt}[1]{\overrightarrow{\mathrm{#1}}}
\providecommand{\vvt}[1]{\overrightarrow{\mathrm{#1}}}
\providecommand{\np}[2][]{{#2~#1}}
\providecommand{\vv}[1]{{\overrightarrow{#1}}}
\renewcommand{\d}{{\mathrm{d}}}
\providecommand{\degres}{{^\circ}}
\providecommand{\degre}{{^\circ}}
\providecommand{\blue}{\color{blue}}
\providecommand{\red}{\color{red}}
\providecommand{\enskip}{\quad}
\providecommand{\llbracket}{\lbrack\!\lbrack}
\providecommand{\rrbracket}{\rbrack\!\rbrack}
\providecommand{\euro}{{\text{€}}}
\providecommand{\no}{{\text{n°}}}
\providecommand{\interval}[2]{{[#1;#2]}}
\providecommand{\intervFO}[2]{{[#1;#2[}}
\providecommand{\dotfill}{\ldots}

\providecommand{\ds}{\displaystyle}
\providecommand{\textding}[1]{#1}
\providecommand{\Oij}{\(\mathrm{O};\vec{\imath},\vec{\jmath}\)}
\providecommand{\Oijk}{\(\mathrm{O};\vec{\imath},\vec{\jmath},\vec{k}\)}
\providecommand{\Ouv}{\(\mathrm{O};\vec{u},\vec{v}\)}
\providecommand{\intervFF}[2]{\ensuremath{[#1;#2]}}
\renewcommand{\euro}{\ifmmode\text{\texteuro}\else\texteuro\fi}
\DeclareUnicodeCharacter{00B0}{\ifmmode^{\circ}\else\textdegree\fi}
\DeclareUnicodeCharacter{20AC}{\ifmmode\text{\texteuro}\else\texteuro\fi}
% Énoncés : Éducation nationale. Transcription et corrigés : APMEP, Denis Vergès et contributeurs. Découpage et classement : CoopMaths, Sébastien Lozano, Éric Elter et contributeurs. Adaptation de la présentation : Maths Explorer.
% CC BY-SA : https://coopmaths.fr/www/about/
% https://coopmaths.fr/alea/static/bac/2023/tex/bac_2023_03_sujet1_reunion_3_cor.tex
\begin{document}
\noindent\textit{APMEP / CoopMaths. CC BY-SA. Présentation adaptée par Maths Explorer.}\par\medskip
\medskip

On considère la suite $\left(u_n\right)$ définie par $u_0 = 3$ et, pour tout $n$ de $\N$:,
$u_{n+1} = \dfrac12u_n + \dfrac12n + 1$.

\bigskip

\textbf{Partie A}

\medskip





\begin{enumerate}
\item La valeur de $u_2$ est égale à :

\begin{center}
\begin{tabularx}{\linewidth}{*{2}{X}}
\textbf{a.~~} $\dfrac{11}{4}$& \textbf{b.~~} $\dfrac{13}{2}$\\[7pt]
\textbf{c.~~} 3,5&\textbf{d.~~}2,7
\end{tabularx}
\end{center}

\medskip

\begin{tabular}{@{\hspace*{0.05\linewidth}} || p{0.93\linewidth}}
$u_1=u_{0+1} = \dfrac{1}{2}u_0 + \dfrac{1}{2}\times 0 +1 
= \dfrac{5}{2}$\\[7pt]
$u_2=u_{1+1} = \dfrac{1}{2}u_1 + \dfrac{1}{2}\times 1 +1 
= \dfrac{1}{2}\times \dfrac{5}{2} + \dfrac{1}{2} +1
= \dfrac{5}{4}+ \dfrac{2}{4} + \dfrac{4}{4}
= \dfrac{11}{4}$
\hfill\textbf{Réponse a.}
\end{tabular}

\medskip

\item La suite $\left(v_n\right)$ définie, pour tout entier naturel $n$, par $v_n = u_n - n$ est :

\begin{center}
\begin{tabularx}{\linewidth}{*{2}{X}}
\textbf{a.~~} arithmétique de raison $\dfrac12$&\textbf{b.~~} géométrique de raison $\dfrac12$\\
\textbf{c.~~}  constante. &\textbf{d.~~}  ni arithmétique, ni géométrique.
\end{tabularx}
\end{center}

\medskip

\begin{tabular}{@{\hspace*{0.05\linewidth}} || p{0.93\linewidth}}
$v_{n+1}=u_{n+1}-(n+1)=\dfrac{1}{2}u_n+\dfrac{1}{2}n+1-n-1 = \dfrac{1}{2}u_n -\dfrac{1}{2}n = \dfrac{1}{2}\left (u_n-n\right )=\dfrac{1}{2}v_n$\\[7pt]
Donc la suite $(v_n)$ est géométrique de raison $\dfrac{1}{2}$.
\hfill\textbf{Réponse b.}
\end{tabular}

\medskip

\item On considère la fonction ci-dessous, écrite de manière incomplète en langage Python.

\begin{minipage}{0.58\linewidth}
$n$ désigne un entier naturel non nul. 

On rappelle qu'en langage Python
\og i in range (n)\fg{} signifie que i varie de 0 à $n-1$.
\end{minipage}\hfill
\begin{minipage}{0.38\linewidth}
\begin{tabular}{|c|l|}\hline
1&def terme (n)\\
2&\qquad U=3\\
3&\qquad for i in range(n): \\
4& \dotfill \\
5&\qquad return U\\ \hline
\end{tabular}
\end{minipage}

Pour que terme (n) renvoie la valeur de $u_n$, on peut compléter la ligne 4 par : 

\begin{center}
\begin{tabularx}{\linewidth}{*{2}{X}}
\textbf{a.~~}U = U/2 + (i$+$1)/2+1	&\textbf{b.~~} U = U/2 + n/2 + 1\\
\textbf{c.~~}U = U/2 + (i$-$1)/2+1 	&\textbf{d.~~} U = U/2 + i/2 + 1
\end{tabularx}
\end{center}

\begin{tabular}{@{\hspace*{0.05\linewidth}} || p{0.93\linewidth}}
Pour $i=0$, il faut calculer $u_1$; la seule formule valable est: U = U/2 + i/2 + 1.
\hfill\textbf{Réponse d.}
\end{tabular}

\end{enumerate}

\bigskip

\textbf{Partie B}

\medskip

\begin{enumerate}
\item On va démontrer par récurrence pour tout entier naturel $n$ la propriété $\mathcal{P}_n$: $n \leqslant u_n \leqslant n + 3$.

\begin{list}{\textbullet}{}
\item \textbf{Initialisation}

Pour $n=0$, on a $u_0=3$ donc $0\leqslant u_0 \leqslant 0+3$.\\
La propriété est vraie au rang 0.
\item \textbf{Hérédité}

Supposons la propriété vraie au rang $n$, c'est-à-dire $n \leqslant u_n \leqslant n + 3$; c'est l'hypothèse de récurrence.

$\begin{gathered}n \leqslant u_n \leqslant n + 3
\iff \dfrac{1}{2}n \leqslant \dfrac{1}{2} u_n \leqslant \dfrac{1}{2}n + \dfrac{3}{2}
\iff \dfrac{1}{2}n +\dfrac{1}{2}n\leqslant \dfrac{1}{2} u_n + \dfrac{1}{2}n\leqslant \dfrac{1}{2}n + \dfrac{3}{2} + \dfrac{1}{2}n\\[7pt]
\phantom{n \leqslant u_n \leqslant n + 3}
\iff n\leqslant \dfrac{1}{2} u_n +\dfrac{1}{2}n\leqslant n + \dfrac{3}{2}
\iff n+1 \leqslant \dfrac{1}{2} u_n \dfrac{1}{2}n +1\leqslant n + \dfrac{3}{2}+1\\[7pt]
\phantom{n \leqslant u_n \leqslant n + 3}
\iff n+1 \leqslant u_{n+1}\leqslant n + \dfrac{5}{2}\end{gathered}$

Or $\dfrac{5}{2}\leqslant 4$, donc on a: $n+1 \leqslant u_{n+1} \leqslant (n+1)+3$.

Donc la propriété est vraie au rang $n+1$.
\item \textbf{Conclusion}

La propriété est vraie au rang 0, elle est héréditaire pour tout $n\geqslant 0$, donc, d'après le principe de récurrence, elle est vraie pour tout $n$.
\end{list}

On a donc démontré que, pour tout entier naturel $n$, on a: $n\leqslant u_n \leqslant n+3$.

\item 
$\ds\lim_{n\to +\infty}n = +\infty$; or pour tout $n$, on a $n\leqslant u_n$. On en déduit que $\ds\lim_{n\to +\infty} u_n = +\infty$

\item 
Pour tout $n$, on a $n \leqslant u_n \leqslant n+3$ donc, pour tout $n$ non nul, on a 
$1 \leqslant \dfrac{u_n}{n} \leqslant 1+\dfrac{3}{n}$.

$\ds\lim_{n\to +\infty} \dfrac{3}{n}=0$ donc $\ds\lim_{n\to +\infty} 1+\dfrac{3}{n}=1$, et $\ds\lim_{n\to +\infty} 1=1$.

Donc, d'après le théorème des gendarmes, on peut dire que la suite $\left(\dfrac{u_n}{n}\right)_{n\in\N^*}$ converge vers 1.

\end{enumerate}

\bigskip

\end{document}
